% protection_layers.tex — Arquitectura defensiva 5 capas (v0.18.0)
\begin{table}[htbp]
\centering
\caption{Arquitectura de protección de datos --- 5 capas}
\label{tab:protection-layers}
\begin{tabular}{@{}clp{8cm}@{}}
\toprule
\textbf{Capa} & \textbf{Componente} & \textbf{Protección} \\
\midrule
L1 & \texttt{meter\_read\_all()} & \texttt{memset(0)} inicial + \texttt{field\_mask} por OBIS individual. Campos fallidos permanecen en 0 con bit OFF. \\
\addlinespace
L2 & \texttt{MIN\_READ\_PERCENT} & Si cobertura de lectura $<50\%$ del objetivo, \texttt{valid = false} y se descarta toda la lectura. \\
\addlinespace
L3 & \texttt{readings\_sanity\_check()} & Rangos físicos: $V \in [50, 500]$\,V, $f \in [40, 70]$\,Hz. Requiere al menos tensión o frecuencia presente. \\
\addlinespace
L4 & \texttt{PUSH\_FIELD(bit\_idx)} & Omite campos cuyo bit no está activo en \texttt{field\_mask}. Solo datos realmente leídos llegan al servidor. \\
\addlinespace
L5 & \texttt{consecutive\_meter\_failures} & Log crítico tras 5 fallos consecutivos de comunicación RS-485/DLMS. \\
\bottomrule
\end{tabular}
\end{table}
